\documentclass[11pt,a4paper]{article}
\usepackage[paperwidth=210mm,paperheight=700mm,margin=1.5cm]{geometry}
\usepackage{graphicx}
\usepackage{float}
\usepackage{caption}
\usepackage{amsmath}
\usepackage[hidelinks]{hyperref}
\usepackage{xcolor}

% No page numbers
\pagestyle{empty}
\makeatletter
\let\ps@plain\ps@empty
\makeatother

\setlength{\parindent}{0pt}
\setlength{\parskip}{0.5em}
\captionsetup{font=small,labelfont=bf}

% Change this to "results" if this tex file is compiled from the exercise folder itself.
\newcommand{\plotdir}{../results}

\newcommand{\resultpair}[3]{%
\begin{figure}[H]
    \centering
    \begin{minipage}[t]{0.48\linewidth}
        \centering
        \includegraphics[width=\linewidth]{\plotdir/#1}
    \end{minipage}
    \hfill
    \begin{minipage}[t]{0.48\linewidth}
        \centering
        \includegraphics[width=\linewidth]{\plotdir/#2}
    \end{minipage}
    \caption{#3}
\end{figure}
}

\title{Transformer Exercise: Neutrino Position Reconstruction}
\author{}
\date{}

\begin{document}

\maketitle

This is a short summary of my transformer model for the 2D IceCube-inspired neutrino dataset. On the test set I got an MSE loss of 2.0061, an MAE of 0.8733, and a mean position error of 1.3731. The mean position error is the average distance between the predicted and true interaction positions.

Each event is a variable-length list of photon hits. Each hit has three inputs: detection time, detector $x$ position, and detector $y$ position. I embedded these features, passed them through a transformer encoder, used a mask to ignore padding, mean-pooled the real hits, and predicted the interaction position $(x,y)$.

The transformer uses 4 attention heads and 2 encoder layers, so it satisfies the exercise requirement. The padding and masking part is needed because different events have different numbers of photon hits.

\section*{Training Curve}

\begin{figure}[H]
    \centering
    \includegraphics[width=0.78\linewidth]{\plotdir/loss_curve.png}
    \caption{Training and validation MSE over epochs.}
    \label{fig:loss-curve}
\end{figure}

\section*{Overall Position Reconstruction}

\begin{figure}[H]
    \centering
    \includegraphics[width=0.70\linewidth]{\plotdir/test_predictions.png}
    \caption{True and predicted neutrino positions in the $(x,y)$ plane. The predictions cover the same general detector area as the true positions, but they are more spread out and less clean near the edges.}
    \label{fig:test-predictions}
\end{figure}

\section*{Position Error Distribution}

\begin{figure}[H]
    \centering
    \includegraphics[width=0.72\linewidth]{\plotdir/position_error_histogram.png}
    \caption{Distribution of test position errors. Most events have a fairly small error, but there is a tail of harder events with much worse reconstruction.}
    \label{fig:position-error-histogram}
\end{figure}

\section*{Coordinate-wise Prediction Quality}

\resultpair
    {pred_vs_true_x.png}
    {pred_vs_true_y.png}
    {The dashed line is perfect prediction. The points roughly follow the line, so the model is learning the position, but the scatter shows the remaining uncertainty.}

\section*{Residual Diagnostics}

\resultpair
    {x_bias_vs_radius.png}
    {y_bias_vs_radius.png}
    {The residuals get wider toward the detector edge, so edge events seem harder to reconstruct. There is no obvious huge global offset.}

\end{document}
